\chapter{West Nile Virus Antibodies}

\section*{What Is This Test?}
West Nile virus (WNV) antibody testing detects antibodies produced after infection with WNV, a mosquito-borne flavivirus. The test is used mainly when WNV encephalitis, meningitis, or acute flaccid paralysis is suspected, and for selected cases of unexplained febrile illness after mosquito exposure.

\section*{Alternative Names}
West Nile virus serology; WNV IgM; WNV IgG; West Nile fever antibody test

\section*{Principle}
Enzyme immunoassays detect WNV-specific IgM and IgG in serum or cerebrospinal fluid. IgM generally indicates recent infection, while IgG develops later and may persist for years. Because flavivirus antibodies cross-react, positive screening results may require confirmatory plaque-reduction neutralization testing.

\section*{Why This Test Is Ordered}
\begin{itemize}
  \item Suspected West Nile encephalitis or meningitis with fever, headache, confusion, neck stiffness, tremor, weakness, or other neurologic findings
  \item Suspected West Nile-associated acute flaccid paralysis
  \item Unexplained febrile illness with compatible mosquito exposure during a period of local transmission
  \item Evaluation of a compatible illness in an immunocompromised patient, who may have an atypical antibody response
\end{itemize}

\section*{Primary vs Secondary Test}
For neuroinvasive disease, WNV IgM in CSF is the preferred diagnostic test. Serum IgM is supportive. WNV PCR is an alternative in selected immunocompromised patients, but a negative PCR does not exclude infection because viremia is often brief and low-level.

\section*{How This Test Is Performed}
Blood is collected for serum testing. When meningitis or encephalitis is suspected, cerebrospinal fluid is tested alongside serum; CSF collection is performed clinically and is not part of the laboratory assay.

\section*{What Preparation Is Needed}
No fasting is required. Provide the date of symptom onset, travel or mosquito exposure, vaccination history, and any prior flavivirus infection when known.

\section*{Risks}
Venipuncture only for serum testing. Risks of CSF collection relate to the collection procedure, not the antibody assay.

\section*{Understanding Results}
\begin{itemize}
  \item \textbf{Positive WNV IgM in CSF:} Strongly supports recent neuroinvasive WNV infection in a compatible clinical setting.
  \item \textbf{Positive serum IgM:} Supports recent infection but may persist for months and can cross-react with other flaviviruses.
  \item \textbf{Positive IgG alone:} Usually indicates previous exposure and does not prove that the current illness is caused by WNV.
  \item \textbf{Negative early testing:} May require repeat serum testing if collected before antibodies develop or if clinical suspicion remains high.
\end{itemize}

\section*{Follow-up Tests}
\begin{itemize}
  \item Cerebrospinal-fluid cell count, protein, and glucose when CNS infection is suspected
  \item Confirmatory plaque-reduction neutralization testing for equivocal or cross-reactive flavivirus results
  \item Testing for other causes of encephalitis or meningitis according to the clinical presentation
\end{itemize}

\clearpage
\section*{Regulatory Status and Authorization Date}
This chapter describes a test category, not a specific commercial product. FDA, CDSCO, EU IVDR, and MHRA approval or registration dates therefore cannot be assigned without the assay manufacturer, product name, instrument, and intended use. For laboratory-developed tests, an individual agency approval date may not exist. See the Preface for the interpretation of regulatory status.

\clearpage
